\section{Method}
\label{sec:method}

%\subsection{Overview}
The overview of our method is illustrated in Fig.~\ref{fig:framework}:
%We are given a pre-trained large language model $M$ and a dataset $\Dcal=\{x_i, y_i\}_{i=1}^D$, where $x_i$ is a question and $y_i$ is the corresponding answer. 
%We use few-shot Chain-of-Thought (CoT) examples~\citep{Wei2022ChainOT} that contain questions, reasoning, and answers to prompt $M$ for generating reasoning paths and answers of questions in the training set $\Dcal^{\mathtt{train}}$. 
%To improve the reasoning accuracy, we apply multiple path decoding with a sampling temperature $T>0$, and use majority voting to select the most consistent final answer, which has shown to be much more effective than greedy decoding~\citep{Wang2022SelfConsistencyIC}. We then filter the reasoning paths that lead to the most consistent answer as the self-training data, and design mixed formats of prompts and answers for augmentation. At testing stage, we will use the self-trained model to do multiple path decoding again to obtain the final answers.
We are given a pre-trained Large Language Model (LLM) $M$ and a question-only training dataset $\Dcal^{\mathtt{train}}=\{x_i\}_{i=1}^D$ with few-shot Chain-of-Thought (CoT) examples~\citep{Wei2022ChainOT}.
%, where $x_i$ is a question and $y_i$ is the corresponding answer. 
%We use few-shot Chain-of-Thought (CoT) examples~\citep{Wei2022ChainOT} that contain questions, reasoning, and answers to 
We apply multiple path decoding with a sampling temperature $T>0$ for generating $m$ reasoning paths and answers $\{r_{i_1}, r_{i_2}, \dots, r_{i_m}\}$ for each question $x_i$ in $\Dcal^{ \mathtt{train}}$, and use majority voting (self-consistency) to select the most consistent, highest confidence answer~\citep{Wang2022SelfConsistencyIC}. We then keep all reasoning paths that lead to the most consistent answer,  apply mixed formats of prompts and answers for augmentation, and fine-tune the model on these self-generated reasoning-answer data. We consider our approach as making the model self-improve. %At test time, we will use the self-finetuned model to with multiple path decoding again to obtain the final answers.
% $\Dcal=\{\Dcal^{\mathtt{train}}, \Dcal^{\mathtt{test}}\}$
In the following sections, we detail important designs within our method, along with additional approaches for the model to self-improve without supervised data.

% \subsection{Generating Multiple Reasoning Paths for Training Questions}
\subsection{Generating and Filtering Multiple Reasoning Paths}% for Training Questions}
%Unlike standard prompting~\citep{Brown2020LanguageMA} that directly pairs questions and answers as few-shot exemplars, Chain-of-Thought prompting~\citep{Wei2022ChainOT} associate answers with text sequences to mimic human's reasoning procedure, and has shown to be significantly more effective in prompting language to do multi-step reasoning tasks. Recently, \textit{Self-Consistency}~\citep{Wang2022SelfConsistencyIC} is proposed to replace greedy decoding with sampling diverse reasoning paths, and use majority voting to select the most consistent answer. 
Self-consistency~\citep{Wang2022SelfConsistencyIC} brings large improvements on reasoning tasks (e.g., $56.5\%\rightarrow74.4\%$ on GSM8K test set), and the gap between greedy decoding and diverse decoding shows there is a potential for further improving the reasoning ability of $M$, using the self-selected high-confidence reasoning paths as training data.
\newcommand{\wronglabel}[1]{\textcolor{red}{(\textit{#1})}}
\begin{table}[h]
\caption{Examples of 3 self-generated CoT reasoning paths given a question. Output 1 and 3 are the most consistent reasoning paths based on majority voting and kept as self-training data.}
% \vspace{-1.5em}
\centering
% \small
\resizebox{\textwidth}{!}{
\begin{tabular}{l}
\toprule
\textbf{Question}: Stefan goes to a restaurant with his family. They order an appetizer that costs \$10 and\\ 4 entrees that are \$20 each. If they tip 20\% of the total, what is the total amount of money that\\ they spend?\\
\midrule
\textbf{Multiple Path Decoding}:\\
\textbf{Output 1}: The appetizer costs \$10. The entrees cost \$20 each so in total \textcolor{blue}{4 * \$20 = \$80}. This means\\ the total cost is \textcolor{blue}{\$10 + \$80 = \$90}. They tip 20\% of it, so the total amount they spend is \textcolor{blue}{\$90 * 1.2}\\ \textcolor{blue}{= \$108}. The answer is 108.\\
\\
\textbf{Output 2}: The appetizer costs \$10 and the entrees are \$20 each. There are 4 entrees so the sum is\\ \$20 * 4 = \$80. The waiter gets 20\% of the total. 20\% of \$80 is \$80 * .2 = \$16. The answer is \$80\\ + \$16 = \$96. \wronglabel{Incorrect reasoning path}\\
\\
\textbf{Output 3}: The appetizer costs \$10. The entrees cost \textcolor{blue}{4 * \$20 = \$80}. The tip is 20\% of the total, so\\ it is 20\% of the \$90 they have spent. The tip is \textcolor{blue}{0.2 * 90 = \$18}. The total they spent is \textcolor{blue}{\$90 + \$18}\\ \textcolor{blue}{= \$108}. The answer is 108.\\
\bottomrule
\end{tabular}
}
\label{tab:self_consist}
% \vspace{-1em}
\end{table}



%To generate extensive high quality Chain-of-Thoughts as training data, we utilize the questions in the training set $\{x_i\}_{i=1}^{D\mathtt{train}}$ without their ground truth answers. 
For each training question ${x_i}$, we sample $m$ CoT reasoning paths, denoted as $\{r_{i_1}, r_{i_2}, \dots, r_{i_m}\}$ (see Table~\ref{tab:self_consist} for examples). 
%An example of a training question with the self-generated CoT reasoning paths is shown in Table~\ref{tab:self_consist}.
Since $M$ is prompted with the CoT examples from~\citet{Wei2022ChainOT}, we apply the same output parsing with ``The answer is'' to generate their predicted answers  $\{y_{i_1}, y_{i_2}, \dots, y_{i_m}\}$.
%Since $M$ is prompted with CoT examples that ends with ``The answer is $y$.'', it is also likely to generate reasoning paths for new questions $x_i$ that end with similar patterns.
%Therefore, we can parse from each reasoning path to derive their final answers $\{y_{i_1}, y_{i_2}, \dots, y_{i_m}\}$. 
The most consistent answer, which is not necessarily a correct answer, is selected by majority voting, denoted as $\tilde{y}_i=\argmax_{y_{i_j}}\sum_{k=1}^m\mathbb{I}(y_{i_j}=y_{i_k})$.
%In Table~\ref{tab:self_consist}, the most consistent answer $\tilde{y}$ is $108$, derived by output path $1$ and output path $3$, while the output path $2$ makes a mistake in calculating the cost of the foods.
For all the training questions, we filter the CoT reasoning paths that reach $\tilde{y}$ as the final answer to be put into the self-training data, denoted as $\Dcal^{\mathtt{self-consistent}}=\{x_i, \tilde{\rvec_i}\}$, where $\tilde{\rvec_i} = \{r_{i_j}|1\leq j \leq m, y_{i_j} = \tilde{y}_i\}$.

\begin{wrapfigure}{wr}{0.4\textwidth}
\vspace{-1.5em}
\begin{center}
    \includegraphics[width=\textwidth]{Figures/scatter.pdf}
\end{center}
\vspace{-0.4cm}
\caption{The relation of accuracy and confidence of the majority-voted answer after multiple path decoding on GSM8K training-set questions. Predicted confidence from self-consistency~\citep{Wang2022SelfConsistencyIC} is well calibrated~\citep{guo2017calibration}.}
\vspace{-0.3em}
\label{fig:scatter}
\end{wrapfigure}
   

Since we do not use any ground truth labels to filter out cases where $\tilde{y}_i\neq y_i$, it is important that the self-generated CoT reasoning paths are mostly reliable and incorrect answers do not hurt the self-improvement of the model.
We plot the relation between the accuracy and confidence of self-generated CoT paths for each question in GSM8K training set in Fig.~\ref{fig:scatter}. 
The confidence is the number of CoT paths leading to $\tilde{y}$ divided by the total path number $m$. The y-axis shows the accuracy of $\tilde{y}$ under a certain confidence. The circle area and the color darkness shows the number of questions under a certain confidence. 
We can observe that confident answers are more likely to be correct, which means that when a question has many consistent CoT paths, then the corresponding $\tilde{y}$ is more likely to be correct. On the other hand, when $\tilde{y}$ is wrong, it is likely to be supported by fewer CoT paths, and brings little noise to the training samples.

%We can observe that the predicted confidence is well calibrated~\citep{guo2017calibration}, i.e. more confident answers are more likely to be actually correct. This means that when $\tilde{y}$ is correct, it likely has many consistent CoT paths, and %, then the corresponding is more likely to be correct. 
%On the other hand, 
%when $\tilde{y}$ is wrong, it is likely to be supported by fewer CoT paths and automatically weighs itself down in the fine-tuning process.

% \vspace{+2em}
\subsection{Training with Mixed Formats}\label{sec:mix_formats}
To prevent the language model from overfitting to specific prompts or answer styles, we create four different formats for each reasoning path to be mixed in the self-training data, shown in Table~\ref{tab:mix_format}. In the first format, a few Chain-of-Thought examples (questions followed by reasoning paths leading to the correct final answers) are prepended to the new question, while the language model output is trained to be the same with the filtered CoT reasoning paths. In the second format, we use examples of questions and their direct answers as standard prompting, and the language model output is supposed to also only contain the direct answer. The third and fourth format are similar to the first and second format, except that no example of question-answer pairs are given, so that the model will learn to think on its own in an in-context zero-shot manner. In the third format, where we want the model to output CoT reasoning without prepending examples containing CoT reasonings, we append ``Let's think step by step.'' at the end of the input sequence, to guide the language model to generate step-by-step CoT reasoning paths \citep{step_by_step}. The mixed formats of training samples are then used to fine-tune the pre-trained language model $M$.

\begin{table}[h]
\caption{An example of how a reasoning path is augmented into four formats of training data with different prompts (in input) and answer styles (in output). Specifically, the \textit{CoT prompting examples} used for each tasks are listed in Appendix~\ref{sec:app_prompt}. The \textit{Standard prompting examples} are the same question-answer pairs with \textit{CoT prompting examples}, except that reasoning is removed.}
% \vspace{1.5em}
\centering
% \small
\resizebox{\textwidth}{!}{
\begin{tabular}{l}
\toprule
\textbf{Question}: Amy is 10 years old. Jake is 8 years old. Alex's age is right in the middle. How old is Alex?\\
\textbf{Selected Chain-of-Thought}: Amy is 10 years old. Jake is 8 years old. Alex's age is in the middle of\\ Amy and Jake, so Alex is ( 8 + 10 ) / 2 = 9 years old. The answer is 9.\\
\midrule
\textbf{Mixed-formats of training data}:\\
\textbf{Format 1}: \textbf{Input}: \textit{[CoT prompting examples]} + `\textbackslash n' + \textit{[Question]} +  `\textbackslash n' + `A:'\\
\textbf{Output}: Amy is 10 years old. Jake is 8 years old. Alex's age is in the middle of Amy and Jake, so Alex\\ is ( 8 + 10 ) / 2 = 9 years old. The answer is 9.\\
\\
\textbf{Format 2}: \textbf{Input}: \textit{[Standard prompting examples]} + `\textbackslash n' + \textit{[Question]} +  `\textbackslash n' + `A:'\\
\textbf{Output}: The answer is 9.\\
\\
\textbf{Format 3}: \textbf{Input}: \textit{[Question]} +  `\textbackslash n' + `A: Let's think step by step.'\\
\textbf{Output}: Amy is 10 years old. Jake is 8 years old. Alex's age is in the middle of Amy and Jake, so Alex\\ is ( 8 + 10 ) / 2 = 9 years old. The answer is 9.\\
\\
\textbf{Format 4}: \textbf{Input}: \textit{[Question]} +  `\textbackslash n' + `A:'\\
\textbf{Output}: The answer is 9.\\
\bottomrule
\end{tabular}
}
\label{tab:mix_format}
% \vspace{-1em}
\end{table}


\subsection{Generating Questions and Prompts}\label{sec:q_gen}
% \subsection{Other Approaches for Self-Improvement}\label{sec:q_gen}
Given a set of training questions and a few human-written Chain-of-Thought (CoT) examples as prompts, our proposed approach enables model self-improvement. However, when the amount of training questions or CoT examples is limited, our method may not generate sufficient training samples for language model self-training. Collecting questions from the web requires human engineering. 
To further reduce human effort, we investigate how to self-generate more training questions as well as example prompts. 

\paragraph{Question Generation}\label{sec:q_gen}
Previous work~\citep{Yoo2021GPT3MixLL, Meng2022GeneratingTD} discuss few-shot data augmentation by generating diverse training samples using LLMs. However, those methods are designed for classification tasks and require ground truth label for each few-shot example. We use a simple yet effective approach to generate diverse questions (without ground truth answers) for in-domain questions. Specifically, we randomly select several existing questions, concatenate them in a random order as input prompt, and let the language model generate consecutive sequences as new questions. We repeat the process to obtain a large set of new questions, then use self-consistency~\citep{Wang2022SelfConsistencyIC} to only keep the questions that have a highly confident answer. Those questions are then used as self-generated training questions.

\paragraph{Prompt Generation}
Given a set of questions, humans can write CoT examples as reasoning paths leading to the final answer. In zero-shot setting without manual prompts, we can generate these CoT paths using the model itself. Following~\cite{step_by_step}, we start the answer with ``A: Let's think step by step.'' and let the language model generate the consecutive reasoning paths. We then use those generated reasoning paths as examples for few-shot CoT prompting.